\subsection{旋转与层级编码}

JHQ 首先对数据库向量和查询应用方形正交变换 $\Pi\in\mathbb R^{d\times d}$ \cite{jhq}。记 $y=\Pi x$、$q'=\Pi q$，则在实数运算下，
\begin{equation}
 \begin{aligned}
 \Pi^\top\Pi=I
 &\Longrightarrow
 \|\Pi(q-x)\|^2\\
 &=(q-x)^\top\Pi^\top\Pi(q-x)=\|q-x\|^2.
 \end{aligned}
 \label{eq:2-isometry}
\end{equation}
因此，量化之前的旋转保持维度和距离。有限精度矩阵乘法在第 4 章单独评估。

主码给出 $y$ 的粗略重构 $\widehat p(x)$。在评测配置中，每八个坐标组成的子空间使用一个 byte：每个坐标取 $-a$ 或 $+a$ 两档，因此该 byte 记录八个符号位。一般而言，$M$ 个子空间各有 $D_s=d/M$ 维，用 $B$ 位从 $L=2^{B/D_s}$ 个标量等级的 $K=2^B$ 种组合中选择。对 $L$ 为整数的支持参数，JHQ 使用 Cartesian 积
\begin{equation}
 a_i=\sigma\Phi^{-1}\!\left(\frac{i-1/2}{L}\right),\quad
 i=1,\ldots,L,\qquad
 \mathcal C_m=\{a_1,\ldots,a_L\}^{D_s}.
 \label{eq:2-cartesian}
\end{equation}
其中 $\Phi$ 为标准正态累积分布函数，$\sigma$ 为根据向量能量估计的尺度。当 $B=D_s=8$ 时，$L=2$、$\mathcal C_m=\{-a,+a\}^8$，且 $a\simeq0.67449\sigma$。单位归一化输入使尺度接近 $1/\sqrt d$。变换与层级表示均继承自 JHQ。

残差编码用 $\widehat r$ 近似剩余误差 $r=y-\widehat p$，形成重构 $\widehat y=\widehat p+\widehat r$。评测中的残差量化器学习每个子空间内各坐标共享的标量等级，每个坐标使用一个 byte。因此，每个向量在元数据之外存储 $d/8$ 个主码 byte 和 $d$ 个残差 byte。降低精排预算减少残差读取，已存储的编码保持不变。

该残差相对于主码重构计算。相比之下，使用残差编码的 IVF-PQ 减去向量 $x$ 所属列表的质心 $\mu_{\ell(x)}$：
\[
 r_{\rm IVF\text{-}PQ}=x-\mu_{\ell(x)},\qquad
 r_{\rm JHQ}=\Pi x-\widehat p(x).
\]
JHQ 使用 IVF 路由，重构向量时不加回 IVF 质心。

\subsection{主码与精排评分}

主码粗筛比较 $q'$ 与粗略重构：
\[
 \primarydist=\|q'-\widehat p\|^2.
\]
只有短名单中的候选才使用主码和残差码重新评分。实现中的精排分数为
\begin{equation}
 \refineddist=\|q'-\widehat p-\widehat r\|^2+\|q'\|^2.
 \label{eq:2-refined}
\end{equation}
附加项 $\|q'\|^2$ 对同一查询的所有候选都是常数，因此该分数按重构向量的平方距离排序；这一顺序未必与 GT 使用的原始向量顺序相同。分数展开见\cref{eq:5-refined-constant}。

\par\medskip\noindent\begin{minipage}{\linewidth}
\centering\captionsetup{hypcap=false}
\resizebox{\linewidth}{!}{\begin{tikzpicture}[
 font=\small\linespread{1}\selectfont,
 box/.style={draw=ThesisNavy!65,rounded corners=2pt,align=center,
 text width=2.65cm,minimum height=1.8cm,inner sep=4pt},
 arr/.style={-{Latex[length=2mm]},draw=ThesisNavy!80,thick}]
\node[box] (x) at (0,0) {Database $\mathcal X$\\$N$ vectors};
\node[box] (u) at (3.6,0) {IVF lists\\$U(q)$: $n_c$ vectors};
\node[box,fill=ThesisNavy!7] (c) at (7.2,0) {Primary top-$c$\\$C_\alpha(q)$\\up to $c$ vectors};
\node[box,fill=ThesisNavy!7] (r) at (10.8,0) {Refined top-$k$\\$\mathcal R_\alpha(q)$};
\draw[arr] (x) -- node[above,font=\footnotesize] {IVF} (u);
\draw[arr] (u) -- node[above,font=\footnotesize] {$D_{\rm P}$} (c);
\draw[arr] (c) -- node[above,font=\footnotesize] {$D_{\rm H}$} (r);
\node[align=center,text width=6.4cm] at (2.6,-2.0)
 {Query $q'=\Pi q$ stays FP32.\\Primary codes score the broad scan.};
\node[align=center,text width=5.8cm] at (9.7,-2.0)
 {Residual $r=\Pi x-\widehat p$.\\Survivors decode $\widehat p+\widehat r$.};
\end{tikzpicture}
}
\captionof{figure}[候选与评分级联]{JHQ 的候选与评分级联。IVF 选择列表，主码评分形成短名单，残差码精排该短名单，再执行最终 top-$k$ 选择。}
\label{fig:2-cascade}
\end{minipage}\par\medskip

若路由得到 $n_c$ 个向量、主码有 $M$ 个子空间、保留至多 $c$ 个候选，评分工作量为
\[
 W_{\rm primary}=O(n_cM)\ \text{lookups},\qquad
 W_{\rm refine}=O(cd)\ \text{coordinates}.
\]
内存访问和候选选择也会增加搜索时间。

\subsection{建模假设}

随机旋转可以分散坐标能量，为上述高斯启发的等级提供动机。距离保持并不意味着坐标独立且服从零均值高斯分布：固定变换将数据库均值 $\mu$ 映为 $\Pi\mu$，且 $\|\Pi\mu\|=\|\mu\|$。第 8 章将中心化作为另一种编码选择进行考察。

主码等级取自高斯分位区间的中点，而非使重构均方误差（mean-squared error，MSE）最小的等级。对 $X\sim\mathcal N(0,\sigma^2)$ 和零判定阈值，后者为
\[
 \begin{aligned}
 a_{\rm MSE}
 &=\arg\min_a\mathbb E[(X-a)^2\mid X>0]\\
 &=\mathbb E[X\mid X>0]
 =\sigma\sqrt{2/\pi}\simeq0.79788\sigma.
 \end{aligned}
\]
这些数值服务于不同目标。对固定的符号码，任意正的公共主码尺度保持主码排序（\cref{eq:4-sign-ranking}），但残差重构仍取决于尺度。因此，这些代数评分性质不能证明量化最优或原始向量上的召回率。
